\bookchapter{الفصل 26: حدود تقنية الهاتف الحالية}

\chapterintro{بعد الحديث عن “ما بعد الوفاة” للهواتف القديمة، لنعد نظرنا للحاضر. 2024-2026 فترة انفجار تقني للهواتف: AI من السحابة إلى الجيب، والاتصال بالأقمار الصناعية من الرائد للمتوسط، والشاشات القابلة للطي من “تجربة” إلى “استخدام يومي”، وتكنولوجيا البطارية شهدت اختراقاً على مستوى المواد لم يُرَ منذ سنوات. هذا الفصل يراجع هذه الحدود التقنية، لمساعدتك في الحكم “هل تستبدل الآن” و“بماذا تستبدل”.}

\booksection{واحد: هاتف AI: حيلة تسويق أم تغيير حقيقي؟}

منذ 2024،تقريباً جميع شركات Android تطلق هواتفها الرائدة تحت راية “هاتف AI”، وأطلقت Apple أيضاً Apple Intelligence في iOS 18. لكن ما الذي يجعل الهاتف “هاتف AI” حقاً؟ شكلت الصناعة ثلاثة معايير للحكم:

\begin{itemize}
  \item \textbf{النماذج الكبيرة على الجهاز}: قدرة الهاتف على تشغيل نموذج لغوي بحوالي 7B معاملات محلياً (مثل نسخ خفيفة من Qwen وLlama وDeepSeek)، وإكمال مهام تلخيص النص والترجمة وتوليد النصوص دون الاعتماد على السحابة. طريقة الحكم: إذا كانت وظائف AI تعمل في وضع الطيران فهي نموذج جانبي.
  \item \textbf{قوة حوسبة NPU (TOPS)}: قوة حوسبة NPU (وحدة معالجة الشبكة العصبية) تُقاس بـ TOPS (تريليون عملية في الثانية). رقائق 2024 الرائدة (مثل Snapdragon 8 Gen3 وDimensity 9300) قوتها 40-70 TOPS، والجيل الجديد 2025-2026 يتجاوز 100 TOPS، وهذا هو الحد العتادي لتشغيل النماذج الكبيرة جانبياً.
  \item \textbf{سيناريوهات AI Agent}: ليس فقط "أسئلة وأجوبة" بل تنفيذ سلسلة مهام عبر التطبيقات — مثلاً "رتب صور اجتماع أمس في ألبوم وشاركه مع Zhang San"، يستدعي النظام تلقائياً الألبام والملفات وWeChat وغيرها. هذا هو المعيار الأساسي المميز لـ "هاتف AI" عن "هاتف بمساعد AI".
\end{itemize}

\booksection{اثنان: مقارنة استراتيجيات AI للشركات: من يجدّ ومن يلصق العلامات؟}

تختلف مسارات AI بشكل واضح، يمكن التركيز عليها عند الشراء:

\begin{itemize}
  \item \textbf{Apple Apple Intelligence (iOS 18+)}: تركّز على تكامل التطبيقات على مستوى النظام — تلخيص البريد وأولويات الإشعارات وتحسين الكتابة وتحرير الصور بالـ AI. نموذج جانبي أساسي، والمهام المعقدة على السحابة عبر Private Cloud Compute. الميزة هي حماية الخصوصية وسلاسة النظام، والعيب هو بطء إطلاق السيناريوهات الصينية، بعض الوظائف تتاحة تدريجياً في 2025.
  \item \textbf{Google Gemini Nano}: مدمج في سلسلة Pixel 8/9، يدعم "Circle to Search" وMagic Compose وملخص التسجيلات وغيرها. Gemini Nano نموذج صغير جانبي حقيقي، وبيئة مفتوحة المصدر ودودة، وهو التنفيذ المرجعي لهواتف Android AI.
  \item \textbf{نموذج Pangu الكبير من هواوي}: تكامل مع HarmonyOS Xiaoyi، يدعم توليد النصوص جانبياً وتحرير الصور بالـ AI وإزالة العناصر بالـ AI (إزالة المارة) وملخص المكالمات. بعد HarmonyOS NEXT 5.0 تكامل أعمق على مستوى النظام، دمج قدرات AI مع النظام.
  \item \textbf{Xiaomi MiLM}: مدمج في HyperOS، يدعم صور AI وملخص AI ورسام AI والدقة الفائقة. Xiaomi تسلك طريق "التعاون الجانبي السحابي"، تشغيل نماذج خفيفة جانبياً والمهام الثقيلة على السحابة.
  \item \textbf{OPPO (Andes) وvivo (BlueLM)}: تكامل عميق في ColorOS 15 وOriginOS 5، يركز على إزالة العناصر بالـ AI وملخص المستندات بالـ AI وأسئلة الصور بالـ AI والبحث عن الصور باللغة الطبيعية ("ابحث عن صور التقطها العام الماضي على البحر").
\end{itemize}

\textbf{نصيحة الشراء}: إذا كنت تحتاج غالباً AI لمعالجة مستندات العمل وترتيب المعلومات، اختر أجهزة بقوة جانبية قوية وتكامل AI على مستوى النظام؛ إذا كنت تستخدم AI لتحرير الصور أحياناً، الهواتف المتوسطة تكفي.

\booksection{ثلاثة: الاتصال بالأقمار الصناعية: من “حيلة” الرائد إلى معيار المتوسط}

الاتصال بالأقمار الصناعية من أكثر التقنيات نزولاً فعلياً في العامين الماضيين. في 2022 Huawei Mate 60 وiPhone 14 كانت حكراً على الرائد، وفي 2025 نزلت لفئة 2000-3000 يوان.

\begin{itemize}
  \item \textbf{نزول رسائل BeiDuo القصيرة للمتوسط}: وظيفة الرسائل القصيرة عبر قمر BeiDou (يمكن إرسال رسالة + موقع حتى 1000 حرف صيني) نزلت لموديلات متوسطة من Redmi وrealme وiQOO وغيرها، بدون اشتراك، مجانية. مناسبة للمشي في البرية والرحلات البرية وغيرها.
  \item \textbf{مكالمات Tiantong الفضائية}: سلسلة Huawei Mate 60/70 وسلسلة Honor Magic6 وغيرها تدعم مكالمات قمر Tiantong-1، يمكن الاتصال/استقبال المكالمات الفضائية وإرسال الرسائل. تحتاج اشتراك باقة فضائية من المشغل (حوالي 10-200 يوان/شهر)، مناسبة لمن يدخل المناطق غير المأهولة كثيراً.
  \item \textbf{رسائل iPhone الفضائية (SOS)}: iPhone 14 والأحدث يدعم SOS فضائي، يمكن إرسال رسائل نصية لخدمات الطوارئ ومشاركة الموقع. خدمة مجانية لسنتين، متاح التجديد في بعض المناطق.
  \item \textbf{توصيات عملية}: مستخدمو المدن لا يحتاجونه غالباً، لكن عشاق الأعمال الخارجية ورجال الأعمال المسافرين للمناطق النائية يُنصحون بشدة بأجهزة بوظائف فضائية — قد تنقذ الحياة في اللحظة الحاسمة.
\end{itemize}

\booksection{أربعة: نضوج منظومة الشاشات القابلة للطي: هل تصلح كهاتف رئيسي؟}

انتقلت الشاشات القابلة للطي من “تجربة تقنية” إلى “قابلة للاستخدام كهاتف رئيسي”. للحكم على نضوج الشاشة القابلة للطي، انظر لثلاثة أمور: متانة العتاد وتكييف التطبيقات ومنظومة التفاعل.

\begin{itemize}
  \item \textbf{متانة العتاد}: عمر المفصلة يصل عموماً لأكثر من 500,000 مرة (100 مرة يومياً لمدة 13 سنة)، مفصلة القطرة تقلل تجعد الشاشة، ومقاومة الماء IPX8 أصبحت معياراً. زجاج UTG + فيلم حماية للشاشة الداخلية يجعلها لم تعد "هشة".
  \item \textbf{تكييف التطبيقات}: التطبيقات الرئيسية (WeChat وTaobao وDouyin وBilibili وOffice وWPS) كيّفت للشاشة الكبيرة، لكن حوالي 30\% من التطبيقات الطويلة لم تُكيّف، تظهر ممطوطة أو بحواف سوداء. يُنصح بتجربة تطبيقاتك المعتادة قبل الشراء.
  \item \textbf{تفاعل وضع الإيقاف}: عند الطي 90°-120°، بعض التطبيقات تدعم وضع الإيقاف "الشاشة العلوية تعرض المحتوى والسفلية للتحكم"، مناسب لمكالمات الفيديو والصور الذاتية والمشاهدة. الصور الذاتية بالإيقاف تجربة فريدة للشاشات القابلة للطي.
  \item \textbf{تقسيم الشاشة والمهام المتعددة}: على الشاشة الكبيرة يمكن فتح تطبيقين بتقسيم الشاشة، حتى ثلاثة تطبيقات بنوافذ صغيرة. منطق تقسيم الشاشة في HyperOS وColorOS وOriginOS ناضج، لكن تقسيم الشاشة في iOS (بنمط iPadOS) لم يُفتح على iPhone بعد.
  \item \textbf{السعر}: الطي الأفقي (كتابي) نزل من 10,000 يوان لفئة 5000-7000 يوان، والطي العمودي (صدفي) نزل لفئة 4000-5000 يوان. في 2025 أطلقت بعض الشركات طياً عمودياً بفئة 3000 يوان، ودخلت الشاشات القابلة للطي سوق الجمهور فعلاً.
\end{itemize}

\textbf{نصيحة الشراء}: إذا كنت مستخدماً كثيفاً للقراءة ومعالجة المستندات والفيديو، الشاشة الأفقية تستحق الاعتبار؛ إذا كنت تطلب Portable والسمات الاجتماعية، الشاشة العمودية أنسب. يُنصح بتجربة 30 دقيقة في المتجر قبل القرار.

\booksection{خمسة: 6G وWi-Fi 7: سرعة وحدود الجيل القادم من الاتصال}

رغم أن 6G لم يُطبّق فعلاً، لكن التقنيات ذات الصلة في الطريق:

\begin{itemize}
  \item \textbf{تطور 6G}: يخطط 3GPP لإكمال معايير 6G في 2027، وتشغيل تجاري حوالي 2030. السرعة النظرية القصوى لـ 6G تصل لـ 1 Tbps (50 ضعف 5G)، وزمن الوصول منخفض لـ 0.1ms، ودعم "تكامل الاتصال والإحساس"، يمكن للهاتف حتى العمل كرadar منخفض الدقة. الشركات في مرحلة التحقق من النماذج الأولية، والمنتجات الاستهلاكية ما زالت بعيدة.
  \item \textbf{Wi-Fi 7(802.11be)}: منذ 2024 أصبح معياراً للهواتف الرائدة، السرعة النظرية القصوى 46 Gbps (4 أضعاف Wi-Fi 6)، يدعم عرض نطاق 320MHz و4K-QAM وMLO (نقل متعدد الروابط). عملياً تحت راوتر Wi-Fi 7، تصل سرعة الهاتف الرائد لـ 2-5 Gbps، أسرع 2-3 مرات من Wi-Fi 6E. إذا اشتريت هاتف Wi-Fi 7، تذكر اقتران راوتر Wi-Fi 7 للاستفادة من الأداء الكامل.
  \item \textbf{Bluetooth 5.4 وLE Audio}: الهواتف الجديدة 2024-2026 تدعم Bluetooth 5.4 وLE Audio (صوت منخفض الطاقة)، يمكن توصيل سماعتين للاستماع لنفس الأغنية، وبث Auracast يتيح لعدة أشخاص الاستماع لمصدر واحد في الأماكن العامة.
\end{itemize}

\booksection{ستة: بطارية الأنود الكربوني السيليكوني: البطارية أخيراً تشهد اختراقاً على مستوى المواد}

عمر بطارية الهاتف كان موضع انتقاد لسنوات “شحن مرة يومياً”، ومنذ 2024 أدى انتشار تكنولوجيا الأنود الكربوني السيليكوني لتحسن ملموس.

\begin{itemize}
  \item \textbf{المبدأ التقني}: السعة النظرية للأنود الجرافيت التقليدي حوالي 372 mAh/g، والسعة النظرية للسيليكون تصل لـ 4200 mAh/g (حوالي 10 أضعاف)، لكن السيليكون يتمدد 300\% عند الشحن/التفريغ ويتحلل. الأنود الكربوني السيليكوني عبر بنية مركبة من سيليكون نانو-كربون يضبط مشكلة تمدد السيليكون في نطاق مقبول، ويرفع كثافة طاقة البطارية 15\%-30\%.
  \item \textbf{الأداء الفعلي}: الهواتف الرائدة 2024-2025 (مثل Xiaomi 14 Ultra وHonor Magic 6 Ultimate وvivo X100 Pro) التي تستخدم الأنود الكربوني السيليكوني، تحزم بطاريات 5000-6000 mAh في هيكل بسمك 7-8mm، بزيادة حوالي 800-1000 mAh عن الجيل السابق، وتدوم ليوم ونصف مع الاستخدام المكثف.
  \item \textbf{تكامل الشحن السريع مئة واط}: الأنود الكربوني السيليكوني متوافق مع الشحن السريع، شحن 100W/120W في 17 دقيقة لم يعد خبراً. في 2025 بعض الموديلات تدعم 150W-240W، الشحن لـ 80\% في 10 دقائق.
  \item \textbf{نصيحة الشراء}: إذا كنت مستخدماً مكثفاً، اختر أجهزة ما بعد 2024 بتركيبة "أنود كربوني سيليكوني + شحن مئة واط"، وقلق البطارية سيخف بشكل كبير.
\end{itemize}

\booksection{سبعة: نظرة سريعة على الفصل: جملة واحدة لتذكر هذه الحدود}

\begin{itemize}
  \item ثلاثة معايير لهاتف AI: نموذج كبير جانبي، قوة NPU TOPS، سيناريو AI Agent عبر التطبيقات.
  \item الاتصال الفضائي: رسائل BeiDou القصيرة نزلت للمتوسط، مكالمات Tiantong لا تزال رائدة.
  \item الشاشات القابلة للطي: العتاد ناضج، تكييف التطبيقات 70\%، الإيقاف/تقسيم الشاشة تجربة فريدة.
  \item 6G في الطريق، Wi-Fi 7 هبط، Bluetooth LE Audio بدأ الانتشار.
  \item أنود كربوني سيليكوني + شحن مئة واط: اختراق مزدوج للبطارية والشحن.
\end{itemize}

\begin{tcbpractice}{جرّب: اصنع "قائمة متطلبات تقنية" لهاتفك القادم}
خذ ورقة، صنّف بالثلاث فئات “ضروري/إضافي/غير ضروري”، اذكر احتياجاتك في AI والأقمار الصناعية والشاشات القابلة للطي وWi-Fi 7 والشحن السريع وسعة البطارية. هكذا في المرة القادمة في المتجر، تستطيع الحكم بسرعة على “نقاط البيع” التي يروجها البائع: هل هي احتياج حقيقي أم دفع للبيع. مثلاً:

\begin{itemize}
  \item ضروري: بطارية 5000mAh فأكثر، شحن 80W فأكثر، NFC، مقاومة ماء IP68.
  \item إضافي: ملخص AI جانبي، Wi-Fi 7، رسائل BeiDou القصيرة.
  \item غير ضروري: شاشة قابلة للطي (أعتاد الهاتف المسطح)، مكالمات فضائية (لا أسافر بعيداً)، شحن 240W فائق.
\end{itemize}

\end{tcbpractice}

\begin{tcbsummary}{خلاصة الفصل}
\begin{itemize}
  \item ثلاثة معايير لحكم هاتف AI: نموذج كبير جانبي، NPU TOPS، سيناريو AI Agent
  \item استراتيجيات AI للشركات: Apple تركز على الخصوصية والتكامل، Google جانبية مفتوحة المصدر، Huawei/Xiaomi/OPPO/vivo لكل منها نقاط قوة
  \item الاتصال الفضائي: رسائل BeiDou القصيرة نزلت للمتوسط، مكالمات Tiantong لا تزال رائدة
  \item الشاشات القابلة للطي: العتاد ناضج، تكييف التطبيقات حوالي 70\%، وضع الإيقاف وتقسيم الشاشة تجربة فريدة
  \item معايير 6G في 2027، Wi-Fi 7 منتشر، Bluetooth LE Audio بدأ الهبوط
  \item بطارية الأنود الكربوني السيليكوني جعلت 5000-6000mAh معياراً للرائد، وقلق البطارية خف بشكل كبير
\end{itemize}

\end{tcbsummary}
